\section{增量索引（Mode）过程}

\subsection{新文档的处理流程}

LightRAG 支持动态新增文档，其 Mode（Upsert = Update + Insert）过程如下：

\begin{enumerate}
    \item \textbf{实体关系提取}：基于语言模型从新文档中提取实体（Entity）和关系（Edge）
    \item \textbf{Profiling}：为实体和关系生成详细的文本描述
    \item \textbf{Embedding}：为实体描述、关系描述和 Chunk 文本建立向量索引
    \item \textbf{子图合并}：将新文档的子图合并到全局知识图谱中
\end{enumerate}

\subsection{实体对齐与命名一致性}

\begin{warningbox}{实体命名一致性是关键}
不同 Chunk 中可能用不同方式提到同一个实体。LightRAG 在 Entity Extraction System Prompt 中明确规定：
\begin{itemize}
    \item 实体名称\textbf{大小写不敏感}
    \item 每个有意义的单词要大写（如 "Knowledge Graph" 而非 "knowledge graph"）
    \item 确保一致性命名（相同实体使用相同名称）
\end{itemize}
实体对齐主要依赖语言模型生成时的一致性，而非后处理匹配。
\end{warningbox}

\subsection{描述合并策略：MapReduce}

当同一实体或关系出现在多个 Chunk 中时，每个 Chunk 贡献了该实体/关系的一个\textbf{侧面描述}。合并策略采用 MapReduce 模式：

\begin{enumerate}
    \item \textbf{收集}：将所有 Chunk 中提取到的描述汇总
    \item \textbf{判断长度}：
    \begin{itemize}
        \item 描述较短 $\rightarrow$ 用分隔符直接拼接
        \item 描述较长 $\rightarrow$ 基于语言模型做语义压缩（Summarization）
    \end{itemize}
    \item \textbf{同步更新}：\texttt{src\_id} 也做 concatenation，记录描述来源
\end{enumerate}

\begin{importantbox}{语义压缩 Prompt}
LightRAG 的 \texttt{prompts.py} 中定义了 \texttt{summarize\_entity\_descriptions} 这一 Prompt，专门用于将过长的实体/关系描述压缩为更精练的摘要，同时保留核心信息。
\end{importantbox}

\subsection{合并示例}

以两个 Chunk 为例：
\begin{itemize}
    \item \textbf{Chunk 1}：提取出实体 Alex（描述A）、Bob（描述A）以及关系 Alex→Bob = "同事"
    \item \textbf{Chunk 2}：提取出实体 Alex（描述B）以及关系 Alex→Bob = "朋友"
\end{itemize}

合并后：
\begin{itemize}
    \item Alex 的描述 = 描述A \texttt{<sep>} 描述B（侧面叠加）
    \item Alex→Bob 的关系描述 = "同事" \texttt{<sep>} "朋友"
    \item 如果描述总长度超过阈值，触发语义压缩
\end{itemize}

\subsection{本章小结}

LightRAG 的增量索引通过实体对齐、描述合并（MapReduce）和语义压缩三步，实现了新文档的无缝融入。整个过程高度依赖语言模型的能力。


